\documentclass[onecolumn, draftclsnofoot,10pt, compsoc]{IEEEtran}
\usepackage{graphicx}
\usepackage{url}
\usepackage{setspace}
\usepackage[noadjust]{cite}
\bibliographystyle{IEEEtran}
\usepackage{geometry}
\geometry{textheight=10in, textwidth=7in}

% 1. Fill in these details
\def \CapstoneTeamName{			The Polycules}
\def \CapstoneTeamNumber{		63}
\def \GroupMemberOne{			Jackson Wells}
\def \GroupMemberTwo{			Corynna Park}
\def \GroupMemberThree{			Joshua Lioy}
\def \CapstoneProjectName{		Molecules in 3D?! And in color!? That I can hold in my hand? No way!!!}
\def \CapstoneSponsorCompany{		Oregon State University}
\def \CapstoneSponsorPerson{		Victor Hsu}

% 2. Uncomment the appropriate line below so that the document type works
\def \DocType{			%Problem Statement
				%Requirements Document
				Technology Review
				%Design Document
				%Progress Report
				}
			
\newcommand{\NameSigPair}[1]{\par
\makebox[2.75in][r]{#1} \hfil 	\makebox[3.25in]{\makebox[2.25in]{\hrulefill} \hfill		\makebox[.75in]{\hrulefill}}
\par\vspace{-12pt} \textit{\tiny\noindent
\makebox[2.75in]{} \hfil		\makebox[3.25in]{\makebox[2.25in][r]{Signature} \hfill	\makebox[.75in][r]{Date}}}}
% 3. If the document is not to be signed, uncomment the RENEWcommand below
\renewcommand{\NameSigPair}[1]{#1}

%%%%%%%%%%%%%%%%%%%%%%%%%%%%%%%%%%%%%%%
\begin{document}
\begin{titlepage}
    \pagenumbering{gobble}
    \begin{singlespace}
    	%\includegraphics[height=4cm]{coe_v_spot1}
        \hfill 
        % 4. If you have a logo, use this includegraphics command to put it on the coversheet.
        %\includegraphics[height=4cm]{CompanyLogo}   
        \par\vspace{.2in}
        \centering
        \scshape{
            \huge CS Capstone \DocType \par
            {\large\today}\par
            \vspace{.5in}
            \textbf{\Huge\CapstoneProjectName}\par
            \vfill
            {\large Prepared for}\par
            \Huge \CapstoneSponsorCompany\par
            \vspace{5pt}
            {\Large\NameSigPair{\CapstoneSponsorPerson}\par}
            {\large Prepared by }\par
            Group\CapstoneTeamNumber\par
            % 5. comment out the line below this one if you do not wish to name your team
            \CapstoneTeamName\par 
            \vspace{5pt}
            {\Large
                \NameSigPair{\GroupMemberOne}\par
                \NameSigPair{\GroupMemberTwo}\par
                \NameSigPair{\GroupMemberThree}\par
            }
            \vspace{20pt}
        }
        \begin{abstract}
		This document is intended to describe each team member's role in the prject and to give an in-depth look at the technology to be used in the project.
		Investigation of different technologies preceeding development is extremely important, doing so allows for informed decisions as to what tool or hardware will be the best fit for the project.
		Each piece of the project will be researched, and three potential choices for each piece will be considered.
		Of the three competing technologies, one will be selected for use in the project.
		%more
        \end{abstract}     
    \end{singlespace}
\end{titlepage}
\newpage
\pagenumbering{arabic}
\tableofcontents
% 7. uncomment this (if applicable). Consider adding a page break.
%\listoffigures
%\listoftables
\clearpage

% 8. now you write!
\section{Role}
	My individual role in the Polycules team pertains to the development of the project's ability to handle biological data files, general software development and managing communication between the computer and the 3D printer.
	In this role, I will lay the ground work for the program, to allow the other group memebers to familiarize themselves with the software and get an idea where their piece of development might be placed.
	Along with laying ground work and general program management, I will implement the feature that handles the input of biological data files.
	This feature sits at the core of the project, but is also non-functional without the other essential components.
	I will also be addressing communication between hardware devices, making sure the computer can correctly supply the 3D printer and multi-filament interface with the information that they need.

\section{Programming Language}
	\subsection{Overview}
	The selection of a programming language is a crucial decision.
	Programming languages exist in countless number, each having their own specific task or purpose.
	Some are designed to interact with other programs, others are designed for data processing or visiualization.
	The programming language chosen to be used will dicate how we solve the problem of producing polychromatic objects and what intermediary software tools will be put to use.
	\subsection{Criteria}
	Due to the large variance between programming languages, we must consider our group's background knowledge and preferences so that a unanimously usable language is selected.
	That being said, preference and background knowledge will not be the deciding factor for language selection.
	A mid to high-level scripting language will be ideal for project development. 
	Scripts are intended to handle user input and send said input to other programs for processing, which is essentially the task our project aims to accomplish.
	High-level languages enable us to execute powerful commands without having to get into the hairy details of an algorithm, which can slow development exponentially.
	\subsection{Potential Choices}
		\subsubsection{Python}
		Python is a widely known and respected high-level programming language.
		The documentation for this language is endless, just like the packages available for use.
		Created 26 years ago, Python still has a major place in the world of programming and is used in a plethora of disciplines[].
		A key point of Python's popularity is its readability.
		Python file syntax forces indentation of paticular sections of code, making scripts easier to read and comprehend.
		Availability of packages makes python a cut-throat competitor, allowing users to accomplish dense tasks in dismal ammounts of code.
		\subsubsection{Perl}
		Perl is a conglomeration of programming languages, borrowing essential features from multiple sources.
		Like python, Perl is a high-level scripting language, capable of pipeline management and file manipulation.
		Perl differs from Python in its text handling abilities, focusing more so on modification than unification.
		This language has been around for 29 years, mainly being used in bioinformatic applications and system/network management.
		The documentation and number of packages available for Perl is present in much lower levels than python, as it has not reached the same number of disciplines.
		\subsubsection{Bash}
		Bash is a command language, basic in function and familiar to any individual that has made use of a command-line interface at some point.
		This language can run files containing bash commands, called scripts. 
		Because of its nature, Bash is able to run any program that can be compiled or interpreted on a command-line, which is feature of almost all programs or scripts. 
		% more	
	\subsection{Discussion}
	
	\subsection{Conclusion}

\section{Conversion Software}
	\subsection{Overview}
	Software used to convert cif and pdb files into STL or WRL file types.
	\subsection{Criteria}
	- open source
	- can handle PDB and CIF files
	- can output in STL WRL and OBJ format
	- has visualization step
	\subsection{Potential Choices}
		\subsubsection{Pymol}
			- may not be open source
			- python based
			- command line access
		\subsubsection{Jmol}
			- python based
			- diverse in function
			- command line access
		\subsubsection{Chimera}
			- command line access
			- specialized by experienced individuals	
	\subsection{Discussion}
	\subsection{Conclusion}

\section{Hardware Communication}
	\subsection{Overview}
	- important for replication of project
	- can make or break ability to produce prints
	\subsection{Criteria}
	- effectively transport information
	- minor to no confusion from use
	- easily available
	\subsection{Potential Choices}
		\subsubsection{SD Card}
		- built in functionality
		- not always functional with computers
		- small, easily lost
		\subsubsection{Wireless}
		- ease of use
		- can be disconnected, user must know how to reconnect
		- speed may become an issue with large files
		\subsubsection{USB}
		- unsure if actually present
		- no user interaction
		- fast relay of data
	\subsection{Discussion}
	\subsection{Conclusion}

\end{document}
